\section{Sets, Cartesian Products, and Functions (04-Oct-1994)}

\subsection*{Naive Set Theory}

\begin{definition}
A \textbf{set} is a well-defined collection of distinct objects, called the elements of the set. We denote membership by $x \in A$.
\end{definition}

A set can be defined by listing its elements $A = \{a, b, c\}$ or by a characteristic property:
$$ M = \{x \in A \mid p(x)\} $$
(the set of all elements in $A$ that satisfy property $p$).

\begin{definition}[Set Operations]\leavevmode
Let $A$ and $B$ be two sets. We define the following operations:
\begin{itemize}
    \item \textbf{Union:} $A \cup B = \{x \mid x \in A \text{ or } x \in B\}$
    \item \textbf{Intersection:} $A \cap B = \{x \mid x \in A \text{ and } x \in B\}$
    \item \textbf{Difference:} $A \setminus B = \{x \mid x \in A \text{ and } x \notin B\}$
    \item \textbf{Complement:} If $A \subseteq E$, the complement of $A$ with respect to $E$ is $\complement_E A = E \setminus A$.
\end{itemize}
\end{definition}

\begin{remark}
The set containing no elements is called the \textbf{empty set}, denoted by $\emptyset$. The set of all subsets of $A$ is called the \textbf{power set} of $A$, denoted by $\mathcal{P}(A)$.
\end{remark}

\subsection*{The Cartesian Product}

\begin{definition}
Let $A$ and $B$ be two non-empty sets. An \textbf{ordered pair} is a collection of two elements $(a, b)$ where $a \in A$ and $b \in B$, with the property that $(a, b) = (c, d) \iff a = c$ and $b = d$.
\end{definition}

\begin{definition}
The \textbf{Cartesian product} of two sets $A$ and $B$ is the set of all ordered pairs that can be formed with elements from the two sets:
$$ A \times B = \{(a, b) \mid a \in A \text{ and } b \in B\} $$
\end{definition}

\subsection*{Functions (Maps)}

\begin{definition}
Let $A$ and $B$ be two non-empty sets. A \textbf{function} (or map) $f: A \to B$ is a rule of correspondence that associates to each element $x \in A$ a unique element $y \in B$, denoted by $y = f(x)$. 
\end{definition}

\begin{definition}[Types of functions]\leavevmode
A function $f: A \to B$ is called:
\begin{enumerate}
    \item \textbf{Injective:} if $(\forall) x_1, x_2 \in A, x_1 \neq x_2 \implies f(x_1) \neq f(x_2)$. (Equivalently: $f(x_1) = f(x_2) \implies x_1 = x_2$).
    \item \textbf{Surjective:} if $(\forall) y \in B, (\exists) x \in A$ such that $f(x) = y$. (Equivalently: $\text{Im} f = B$).
    \item \textbf{Bijective:} if $f$ is simultaneously injective and surjective.
\end{enumerate}
\end{definition}

\begin{definition}
Let $f: A \to B$ and $g: B \to C$. The \textbf{composition} of the two functions is a new function $g \circ f: A \to C$ defined by:
$$ (g \circ f)(x) = g(f(x)), \quad (\forall) x \in A $$
\end{definition}

\begin{proposition}
The composition of functions is \textbf{associative}. If $f: A \to B$, $g: B \to C$, and $h: C \to D$, then:
$$ h \circ (g \circ f) = (h \circ g) \circ f $$
\end{proposition}